% Definições básicas do resumo
\begin{adjustwidth}{-1cm}{-1cm}\vspace{\onelineskip}

%% Resumo

\begin{resumo}

O desempenho acadêmico no ensino médio público é fortemente condicionado por fatores sociológicos como pobreza, trabalho juvenil, exclusão digital e distorção idade-série, mas o acesso a microdados individuais de estudantes é restringido por barreiras legais e institucionais. Este trabalho propõe um modelo preditivo de tendência de desempenho acadêmico para estudantes da rede pública estadual do Piauí, fundamentado em fatores sociológicos e treinado sobre uma população sintética representativa do perfil de uma turma média dessa rede, construída pela própria pesquisa. Cada parâmetro de geração é rastreável a estatísticas oficiais e à literatura acadêmica, e o gerador incorpora mecanismos anti-deterministas: a assistência financeira estudantil é derivada de regra legal de elegibilidade, as vulnerabilidades interagem de forma não linear e perfis de resiliência acadêmica são preservados. Os indicadores acadêmicos observáveis correspondem a parciais do primeiro semestre letivo, e a variável-alvo representa a tendência de situação ao final do ano, rotulada pelos critérios de aprovação da Lei de Diretrizes e Bases da Educação Nacional. Sobre a base de três mil registros, três algoritmos de classificação e dois \textit{baselines} de referência foram comparados sob validação cruzada estratificada, com sobreamostragem confinada a cada partição de treino e teste estatístico pareado entre os modelos. O XGBoost obteve o melhor desempenho, com acurácia de 84,11\% (estatisticamente equivalente à \textit{Random Forest} e superior à \textit{Decision Tree} única e aos \textit{baselines}), e nenhum modelo cometeu erros entre as classes extremas. Após os sinais acadêmicos parciais, os atributos mais determinantes foram a distorção idade-série, o recebimento de assistência financeira, o trabalho juvenil e a renda per capita, recuperando o peso relativo atribuído aos fatores sociológicos no fenômeno gerador (evidência de validade metodológica da abordagem, cuja confirmação empírica depende da transposição a dados reais). Conclui-se que a articulação entre dados sintéticos rastreáveis e classificadores de conjunto constitui alternativa metodológica válida para o desenvolvimento de sistemas de alerta precoce em contextos de dados restritos.
\vspace{\onelineskip}
\noindent

\textbf{Palavras-chaves}: desempenho acadêmico; fatores sociológicos; dados sintéticos; aprendizado de máquina; ensino médio público.

\end{resumo}
\end{adjustwidth}
